\section{Related Literature}


\paragraph{Task Interdependence in Task-based Models of Automation.}
Our framework relates to the literature on task-based models of technical change, which generally view tasks as predetermined and technologically independent objects combined through a CES production function \citep{acemoglu2011skills, acemoglu2018, acemoglu2019automation, acemoglu2024tasks}.
In these settings, task-level comparative advantage pins down factor assignment and automation follows relative input efficiencies.
We depart from this literature by modeling production as a sequence of technologically interdependent steps, where AI chaining creates complementarities across adjacent steps that can overturn comparative advantage as the driver of who does what.\footnote{A related perspective on task interdependence in task-based models appears in \cite{trammell2026workflows}, who studies learning spillovers across tasks.
Although both his framework and ours feature interdependent tasks, their focus and margin of analysis differ sharply.
In our setting, task interdependence is an intrinsic technological feature that creates coordination frictions when production is shared across multiple workers.
By contrast, \cite{trammell2026workflows} models interdependence as an endogenous result of economies of scale for the same worker, where performing more tasks raises productivity on subsequent tasks through learning, abstracting from cross-worker coordination frictions within the firm.
}

The chaining feature of the model further connects to, and helps reconcile, competing views about automation in production.
Although much of the AI and labor literature emphasizes task-level substitution between AI and workers, \cite{autor2003skill}, \cite{acemoglu2019automation}, and \cite{bresnahan2002information}, among others, argue that meaningful substitution often occurs at the system level.
Our model shows how step-level automation decisions can aggregate into system-level changes through AI chaining.
Rather than requiring that steps be automated one by one, entire chains can be automated in a single leap, consistent with the perspective emphasized by \cite{bresnahan2002information}.


\paragraph{Measuring AI Exposure and Realized Execution.}
A rapidly growing literature has sought to quantify the impact of new automation technologies on the labor market.
Early work studies the susceptibility of occupations to computerization \citep{frey2017future}, and more recent work develops measures of occupational exposure to AI \citep{brynjolfsson2018can, webb2020impact, felten2021occupational, eloundou2023gpts, tomlinson2025working}, typically by mapping technologies to specific tasks and then aggregating these task-level mappings linearly to construct occupation-level indices.
We depart from this literature by showing that linear aggregation obscures the crucial role of task interdependence.
In our framework, the productivity gains from AI depend not only on how many of an occupation's steps are exposed to AI, but also on where those steps sit relative to one another in the production sequence.
To capture these non-linear interactions, we introduce the concept of ``fragmentation,'' the degree to which AI-suitable steps are dispersed across the workflow.
In this sense, our fragmentation index serves as a bridge between occupation-level exposure measures and the actual patterns of AI execution observed in practice.


\paragraph{Technology Adoption, Complementarities, and O-Ring Dynamics.}
Our work contributes to the literature on friction-laden technology adoption and organizational complementarities.
Empirical work shows that information technologies raise productivity primarily when combined with complementary organizational changes \citep{bresnahan2002information, bloom2016management}.
This aligns with the supermodularity framework of \cite{milgrom1990economics, milgrom1995complementarities}, in which clusters of mutually reinforcing practices drive performance.
We provide a micro-foundation for these complementarities through the lens of the O-ring theory of production \citep{kremer1993oring}: when tasks are complementary because of sequencing, the returns to improving any one step depend on performance elsewhere in the chain.
In our setting AI chaining makes this logic salient, since the payoff to automating a step depends on whether adjacent steps can be executed as part of the same AI block, generating threshold effects and discrete changes in optimal AI adoption and job design.
Our results align with \citet{gans2026oring}, who likewise emphasize task interdependence in O-ring production and show that complementarities can generate ``lumpy'' automation decisions via time/attention reallocation rather than workflow adjacency.

Our framework also provides a micro-foundation for the ``productivity J-curve'' argument \citep{brynjolfsson2021productivity}, in which adoption costs and transitional reorganization precede realized productivity gains.
In our model, the marginal benefit of improving AI technology is non-monotone, with substantial gains emerging only once the technology crosses thresholds that induce discrete reorganizations of work and enable the formation of longer AI chains.
This dynamic is consistent with recent analyses of U.S. Census microdata, which show that AI adoption initially reduces productivity \citep{mcelheran2025rise}, as well as evidence from several other studies \citep{furman2019ai, tambe2012productivity, bonney2024tracking, mcelheran2024adoption}.
